\title{Collaborative Proposal --- ITR/IIS: Cross-modality Heterogeneous
Web Source Integration for High Quality Internet Information Access}

\section{Overview and Significance}
\label{sec-overview}

\newcommand{\jc}[1]{{\bf (** Junghoo ** }{#1}{\bf\ ** Junghoo **)~}}

This project intends to enable Web-based searches to be both 
accurate and efficient by advancing three fundamental technologies:
{\em query formulation}, {\em cross-modality web-page characterization},
and {\em heterogeneous source integration}.

The amount of Internet information grows rapidly, not only text data,
but also multimedia data such as images, video, and music. 
Furthermore, a substantial amount of information is ``hidden''
behind search forms, making it hard for traditional crawlers to
retrieval.  These trends will make traditional search engines 
increasingly ineffective for three reasons.

\begin{enumerate}

\item {\em Ineffective query articulation}. 
 At present, Internet users convey their
 query concepts via {\em keywords}.  However, many complex
 query concepts cannot be articulated by merely a few words.

\item {\em Inadequate cross-modality meta-data integration}.  
 Traditional Web-based systems focus on text
 information while ignoring multimedia contents.
 Text-based only indexing fails to leverage
 multimedia meta data to enrich information representation.
 In addition, text-only information representation
 cannot effectively support the growing need of multi-modality
 searches.  

\item {\em Growing unindexed pages}.
A large number of Web pages are not indexed by search engines and
the number is growing for two reasons.
\begin{itemize}
\item Stale indexing.  Most search engines crawl Web pages periodically 
  (e.g., once a month).  New Web sites created between two updates
  are not indexed, and hence search results might not be
  up-to-date.  As the amount of information continues to
  grow, this problem exacerbates.

\item Static-page-only indexing.   Many search engines index only static 
  pages.  However, there are a huge number of pages in the hidden Web
  \cite{} are {\em dynamically generated} in response to queries via
  search forms.  Statistics shows the percentage of these
  pages is growing, and static indexing misses out a lot of
  information.
\end{itemize}
\end{enumerate}

Consider that CIA wants to identify all information sources 
containing atomic bomb design in the Internet.
(This example was provided by In-Q-Tel, a CIA branch at Palo Alto.)
Unlike that most queries can be satisfied by just a couple of
good results, queries of this kind demand high recall.
To achieve high recall, a search engine
faces three hurdles:
comprehending the query concept completely, 
characterizing the data at hand adequately, 
and integrating with as much useful sources in the Internet 
as possible.
These three hurdles for achieving the performace goal are
precisely caused by the the three problems enumerated above.
This project aims to investigate novel 
techniques in three thrusts to overcome these three hurdles.

\begin{enumerate}
\item {\em Intelligent query concept formulation through active learning}.
 In order to make query formulation easy and complete, it is both
 necessary (for capturing abstract and subjective concepts) and
 desirable (for alleviating difficulties from specifying complex query
 concepts) to build intelligent search engines that can quickly {\em
 learn} users' query concepts through {\em active learning}.  We plan
 to extend our active learning algorithms $\MEGA$ 
 and $\SVMActive$ \cite{LCL01,SVM01},
 and develop new algorithms for
 cross-modality query-concept learning.

\item {\em Cross-modality information characterization and access}.
 We plan to investigate novel data models for multi-modality data
 characterization and meta-data integration.  We believe that a
 cross-modality model that leverages meta-data of different modality
 in an integrated manner will be more effective than treating
 multi-modality meta-data independently.

\item {\em Timely heterogeneous source integration}. 
 In this thrust, we focus on improving quality of data retrieved from
 the Web.  In particular, we investigate how to keep Web indexes
 up-to-date to maximize data freshness; and how to perform seemless
 integration of static and dynamically-generated Web pages.
\end{enumerate}

We have conducted research in learning, information integration, and
keeping Web indexes up-to-date.  The progress we made recently
provides a good foundation for our research team to continue making
more progress. In addition, the Internet crawler and 
image retrieval prototype
\cite{demo-Sigmod01,demo-ACMmm01} that we have implemented allow us
to quickly examine new hypotheses (some are described in
Section~\ref{sec-future-work}), and to solicit input from the research
community and from our partners.  To evaluate how fast and how well
our developed algorithms capture subjective query concepts, can
leverage cross-modality meta data, and can keep information up-to-data
and comprehensive, we have established several testbeds with
colleagues at IBM Almaden Research Center, IBM T.J. Watson Research
Center, and Benchthalon.  We are also participating in efforts
\cite{Bench01} for establishing standard tests \cite{Benchathlon} for
evaluating information access that involves multimedia data.

\subsection*{\underline{Synergisms and Impact}}

Our three research thrusts are both complementary and synergistic:
Our cross-modality paradigm will enrich Web data representation.
The active learning algorithms will enable users to articulate their 
``information need'' more intuitively and effectively, and thereby
make query formulation more personalized and complete.
And a timely and comprehensive data collected from the Web 
will enhance the query results.
The faculty involved in this proposal will all work on more than one thrust.
We have a record of sharing research interests, working together
and joint publications. Our thrusts are a convenient way to
organize our research agenda, but are certainly not meant to partition 
the project.

This project could have substantial
impact upon image retrieval in general, and
upon education and military intelligence in particular.
For improving education,
rapid proliferation of multimedia
content in digital libraries and on the Web underscores the increasing
importance of having effective cross-modality search tools.  
Intelligent query-concept learners will directly or indirectly make
traditional text-based information retrieval easier and more
personalized.  Directly, a text collection can employ an intelligent
learner to better capture users' query concepts.  Indirectly, for
instance, multimedia data can be added to a text collection so that
searches can be conducted through interfaces that contain pictures and
graphics. As a result, even young students who have not learned 
Boolean algebra can use images and graphics to search for stories 
and books.  For mililary intelligence,
the present world situation makes the gathering of intelligence
more essential than ever.  The ability to find
information that fit a query profile
in a massive heterogeneous database offers a strategic
advantage to those who can assemble such information
quickly and accurately.  

